\documentclass[12pt]{article}
\usepackage[utf8]{inputenc}
\usepackage[T1]{fontenc}
\usepackage{hyperref}
\title{Bolyu Language Resources for K-12 Curriculum Development}
\author{Compiled by Claude Sonnet 5.5 (claude-sonnet-5-5) for Andr\'as Kornai}
\date{}
\begin{document}
\maketitle

\section{Introduction}
Bolyu (ISO 639-3 \texttt{ply}, Glottocode boly1239) is an Austroasiatic language of the Pakanic branch, closely related to Bugan. Search results place it in far western Guangxi (Xilin and Longlin counties), on the Guizhou and Yunnan borders. Speaker estimates range from about 500 to 1,400 (2012). It is rarely written, though a Latin-based writing is possible. The search found no school materials, primers, or community organisations. Very little exists, and this report lists only what was found.

\section{Primary Educational Resources}
No primers, textbooks or school programmes for Bolyu were found in the search results.

\section{Linguistic and Reference Materials}
\begin{itemize}
\item Wikipedia overview of the Bolyu language, which cites the main descriptive works: \url{https://en.wikipedia.org/wiki/Bolyu_language}
\item Omniglot page on Bolyu with an alphabet chart: \url{https://omniglot.com/writing/bolyu.htm}
\item Bolyu verb wordlist in the ABVD database: \url{https://abvd.eva.mpg.de/austronesian/language.php?id=960}
\item Bolyu numerals page (Max Planck Institute, EVA): \url{https://lingweb.eva.mpg.de/channumerals/Bolyu.htm}
\end{itemize}
The overview page cites Li Xulian (1999), \emph{A Study of Lai (Bolyu)}, and Jerold Edmondson's English--Bolyu glossary (1995). Their availability was not confirmed, and no full grammar or published dictionary was found.

\section{Religious and Community Resources}
None were found.

\section{Digital and Audio Resources}
No audio, video or software resources were found, apart from the online wordlists above.

\section{Curriculum Development Context}
Bolyu is not known to be taught in schools. A curriculum would first need documentation: fieldwork, a practical orthography, and a lexicon built on the existing glossary and wordlists. Community partners could not be identified from the search results.
\end{document}